\documentclass{beamer}

\input{../helper}

\title{Introdução à Computação \\ arrays e strings}
\author{Alexandre Rademaker}
\date{}

\begin{document}

\begin{frame}
 \maketitle
\end{frame}

\begin{frame}[fragile,allowframebreaks]{Array}

  \lstinputlisting[style=tinyc,caption=scores0.c]{scores0.c}

  Nomes distintos para cada posição na memória! Não necessariamente
  alocados de forma contínua! Podemos simplificar?

  \framebreak

  Isto torna o código mais
  \href{https://en.wikipedia.org/wiki/Code_smell}{elegante}?

  \lstinputlisting[style=tinyc,caption=scores1.c]{scores1.c}

  \framebreak

  \includegraphics[width=.9\textwidth]{scores.png} 

\end{frame}  


\begin{frame}[fragile,allowframebreaks]{caracteres e inteiros}

  \lstinputlisting[style=normalc,caption=hi0.c]{hi0.c}

  \begin{lstlisting}[style=normalbash]
    > clang hi0.c
    > ./a.out
    HI!
    72 73 33
  \end{lstlisting}

  Os \ilcode{char} são armazenados como inteiros (códigos ASCII).
  \ilcode{printf} fez uma conversão \emph{implícita} usando \ilcode{\%i};
  também podemos converter explicitamente:

  \begin{lstlisting}[style=normalc]
    printf("%i %i %i\n", (int) c1, (int) c2, (int) c3);
  \end{lstlisting}

  Mas conversões podem \emph{perder} informação, por exemplo, de
  \ilcode{float} para \ilcode{int}.

\end{frame}

\begin{frame}[fragile,allowframebreaks]{Strings}

  Representamos \emph{strings} como arrays de caracteres (cada um com
  8 bits ou um byte), um tipo que na verdade não é primitivo em C.

  \lstinputlisting[style=normalc,caption=hi1.c]{hi1.c}

  \framebreak

  \begin{lstlisting}[style=normalbash]
    > ./a.out
    HI! vs HI!
    a[0] = H (72)
    a[1] = I (73)
    a[2] = ! (33)
    a[3] =   (0)
  \end{lstlisting}

  \includegraphics[width=.45\textwidth]{chars.png}\hfill
  \includegraphics[width=.45\textwidth]{string.png}

  Em C, strings terminam com o char \ilcode{'\\0'} ou inteiro
  \ilcode{0} (um byte com todos os bits 0), marcando a
  \emph{terminação} da string. Este é o chamado 'NUL'. Sempre
  precisamos de um byte adicional.

  Outros tipos de dados de tamanho \emph{fixo} não precisam de tal
  mecanismo. Mas bibliotecas que oferecem representação de variadas
  precisões numéricas podem usar também tal recurso.

  \framebreak

  \lstinputlisting[style=normalc,caption=hi2.c]{hi2.c}

  \framebreak

  \includegraphics[width=.9\textwidth]{hi-bye.png}

  \framebreak

  Um array de strings ou um array de array de caracteres.

  \lstinputlisting[style=normalc,caption=hi3.c]{hi3.c}

  \framebreak

  Na memória, este array está armazenado da seguinte forma.

  \includegraphics[width=.9\textwidth]{string-matrix.jpg}

\end{frame}


\end{document}

%%% Local Variables:
%%% mode: latex
%%% TeX-master: t
%%% End:
